\documentclass[11pt,a4paper]{article}
\usepackage[margin=1in]{geometry}
\usepackage{amsmath,amssymb,amsfonts,amsthm}
\usepackage{bm}
\usepackage{booktabs}
\usepackage{enumitem}
\usepackage{hyperref}
\usepackage{cleveref}

\newcommand{\R}{\mathbb{R}}
\newcommand{\tr}{\operatorname{tr}}
\newcommand{\diag}{\operatorname{diag}}
\newcommand{\cond}{\operatorname{cond}}
\newcommand{\rank}{\operatorname{rank}}
\newcommand{\expm}{\operatorname{expm}}
\newcommand{\vect}[1]{\bm{#1}}
\newcommand{\mat}[1]{\mathbf{#1}}
\newcommand{\tbot}{\tau_{\mathrm{bot}}}
\newcommand{\ttrans}{\tau_{\mathrm{tr}}}

\theoremstyle{definition}
\newtheorem{proposition}{Proposition}
\newtheorem{remark}{Remark}

\title{Extending the Magnus Integrator to Thick Atmospheres:\\
Mathematical and Computational Challenges}
\author{Technical Report for the Pythonic-DISORT SSA Profile Retrieval Project}
\date{March 2026}

\begin{document}
\maketitle

%======================================================================
\section{Problem Statement}
%======================================================================

The \texttt{pydisort\_magnus} solver discretises the one-dimensional radiative
transfer equation (RTE) for an atmosphere with continuously
$\tau$-varying single-scattering albedo $\omega(\tau)$ and phase function.
After azimuthal Fourier decomposition and angular discretisation with $2N$
Gauss--Legendre quadrature streams, each Fourier mode satisfies a $2N$-dimensional
linear ODE
\begin{equation}\label{eq:ode}
  \frac{d\vect{I}}{d\tau} = \mat{A}(\tau)\,\vect{I}(\tau) + \vect{S}(\tau),
  \qquad \tau\in[0,\tbot],
\end{equation}
subject to boundary conditions at $\tau=0$ (top of atmosphere) and $\tau=\tbot$
(bottom boundary), where $\vect{I}=[\vect{I}^+, \vect{I}^-]^T$ collects the
upwelling ($\vect{I}^+$, $N$ components) and downwelling ($\vect{I}^-$, $N$
components) discrete-ordinate intensities.

The first-order Magnus integrator discretises~\eqref{eq:ode} into $K$ steps of
size $h=\tbot/K$.  At each step the midpoint-rule matrix exponential
\begin{equation}\label{eq:magnus_step}
  \mat{\Phi}_k = \exp\!\bigl(h\,\mat{A}(\tau_{k+1/2})\bigr)
\end{equation}
advances the homogeneous solution, and the particular-solution increment
$\vect{\delta}_k$ from the collimated beam source is obtained from the augmented
matrix exponential.

\textbf{The problem.}  For \emph{thin} atmospheres ($\tbot\lesssim 2$) the
solver works correctly (19/25 tests pass).  For \emph{thick} atmospheres
($\tbot=5$--$32$), the step-by-step SVD accumulation of the propagator
catastrophically fails (6/25 tests fail).  This report analyses why, and
proposes two viable strategies within the Magnus framework.


%======================================================================
\section{The DISORT Coefficient Matrix}\label{sec:A_structure}
%======================================================================

The coefficient matrix has the block structure
\begin{equation}\label{eq:A_block}
  \mat{A} = \begin{pmatrix} -\mat{\alpha} & -\mat{\beta} \\
                              \mat{\beta}  &  \mat{\alpha} \end{pmatrix},
  \qquad
  \mat{\alpha} = \mat{M}^{-1}(\mat{D}^+\mat{W} - \mat{I}_N),
  \quad
  \mat{\beta}  = \mat{M}^{-1}\mat{D}^-\mat{W},
\end{equation}
where $\mat{M}=\diag(\mu_1,\dots,\mu_N)$ (quadrature cosines),
$\mat{W}=\diag(w_1,\dots,w_N)$ (quadrature weights), and
$\mat{D}^\pm$ are the $N\times N$ phase-function kernel matrices evaluated at
$(\mu_i, \pm\mu_j)$ and scaled by $\omega(\tau)$.

\subsection{Eigenvalue structure}
The eigenvalues of $\mat{A}$ come in exact $\pm\lambda$ pairs
(consequence of the block symmetry in~\eqref{eq:A_block}):
\[
  \lambda_1 > \lambda_2 > \cdots > \lambda_N > 0 > -\lambda_N > \cdots > -\lambda_1.
\]
For isotropic scattering with $N=4$ and $\omega=0.5$:
$\lambda_{\max}\approx 13.77$, $\lambda_{\min}\approx 0.97$.
As $\omega\to 1$, the smallest eigenvalue pair $\lambda_N\to 0$ (diffusion
mode), while $\lambda_1$ remains $O(1/\mu_{\min})\approx 14$.

\subsection{Non-normality}\label{sec:nonnormal}
\emph{$\mat{A}$ is not symmetric and not normal.}  The $\mat{M}^{-1}$ factor
in~\eqref{eq:A_block} breaks symmetry: $\mat{A}\neq\mat{A}^T$ and
$\mat{A}\mat{A}^T\neq\mat{A}^T\mat{A}$.  We measured
\begin{equation}
  \frac{\|\mat{A}\mat{A}^T - \mat{A}^T\mat{A}\|_F}{\|\mat{A}\|_F^2}
  \;\approx\; 0.08 \;(\omega=0.2),\quad 0.19 \;(\omega=0.5),\quad 0.34 \;(\omega=0.99).
\end{equation}
The eigenvector matrix $\mat{V}$ has condition number
$\cond(\mat{V})\approx 1.15$--$3.74$ (increasing with $\omega$).  While
these values appear modest, they have profound consequences for the matrix
exponential.


%======================================================================
\section{The Propagator and Its Singular Values}\label{sec:propagator}
%======================================================================

The formal solution of the homogeneous part of~\eqref{eq:ode} over
$[0,\tbot]$ is the \emph{propagator} $\mat{\Phi}(\tbot)=\exp(\mat{A}\tbot)$
(for constant $\mat{A}$; the variable-coefficient case is a time-ordered
exponential).

\subsection{Singular values vs.\ eigenvalues for normal matrices}
If $\mat{A}$ were normal, we would have
$\sigma_j\!\bigl(e^{\mat{A}\tau}\bigr)=e^{\text{Re}(\lambda_j)\tau}$ for
each eigenvalue $\lambda_j$.  In particular, the $N$ modes with positive
$\lambda_j$ give $\sigma_j\gg 1$ (growing) and the $N$ modes with negative
$\lambda_j$ give $\sigma_j\ll 1$ (decaying), with
$\sigma_{\max}/\sigma_{\min}=e^{2\lambda_{\max}\tau}$.

\subsection{Non-normal transient growth}\label{sec:transient}
For non-normal $\mat{A}$, the singular values of $e^{\mat{A}\tau}$ are
\emph{not} simply $e^{\text{Re}(\lambda_j)\tau}$.  The classical
Kreiss matrix theorem~\cite{kreiss1962} bounds the transient growth:
\[
  \|e^{\mat{A}\tau}\|_2 \;\leq\; \cond(\mat{V})^2\,\max_j e^{\text{Re}(\lambda_j)\tau}.
\]
Our diagnostics reveal a more striking phenomenon: for thick atmospheres,
\emph{all} singular values of $e^{\mat{A}\tau}$ eventually become $\gg 1$,
including those that initially decrease.  \Cref{tab:svd_tau5} shows the
8 singular values of $\exp(\mat{A}\cdot\tau)$ for isotropic scattering,
$\omega=0.5$:

\begin{table}[h]\centering
\caption{Singular values of $\exp(\mat{A}\tau)$, isotropic scattering,
$\omega=0.5$, $N=4$ ($2N=8$ streams).}\label{tab:svd_tau5}
\begin{tabular}{r|rrrr|rrrr}
\toprule
$\tau$ & $\sigma_1$ & $\sigma_2$ & $\sigma_3$ & $\sigma_4$
       & $\sigma_5$ & $\sigma_6$ & $\sigma_7$ & $\sigma_8$ \\
\midrule
0.5  & $10^{2.9}$  & 4.0   & 2.0   & 1.6   & 0.61 & 0.51 & 0.25  & $10^{-3.0}$ \\
1.0  & $10^{6.0}$  & 16    & 3.9   & 2.7   & 0.37 & 0.26 & 0.062 & $10^{-6.0}$ \\
1.5  & $10^{9.0}$  & 65    & 7.6   & 4.3   & 0.23 & 0.13 & 0.015 & $10^{-9.1}$ \\
2.0  & $10^{12.0}$ & 259   & 14.9  & 7.0   & 0.14 & 0.067& 0.004 & $\mathbf{4\times 10^{-7}}$ \\
3.0  & $10^{17.9}$ & 4129  & 57    & 18    & 1.42 & 0.091& 0.068 & $\mathbf{0.019}$ \\
5.0  & $10^{29.9}$ & $10^{12}$ & \multicolumn{2}{c|}{$10^{11}$}
     & $10^{11}$ & $10^{10}$ & $10^{10}$ & $\mathbf{4.7\times 10^{7}}$ \\
\bottomrule
\end{tabular}
\end{table}

\textbf{Key observation.}  At $\tau=1.5$, $\sigma_8=7.9\times 10^{-10}$
(minimum).  By $\tau=2.0$, $\sigma_8$ has \emph{rebounded} to
$3.9\times 10^{-7}$; by $\tau=5.0$ it reaches $4.7\times 10^{7}$.  This
``dip-and-rebound'' is a hallmark of non-normal transient growth: the
interference between non-orthogonal eigenvectors causes all singular values
to eventually align with the dominant growing mode.

\begin{proposition}[All singular values grow for thick atmospheres]
For the DISORT coefficient matrix $\mat{A}$ with $\omega>0$ and
$\tau\to\infty$, all singular values of $\exp(\mat{A}\tau)$ grow as
$\sigma_j \sim C_j\,e^{\lambda_{\max}\tau}$ for some constants $C_j>0$.
This is because $e^{\mat{A}\tau} = \mat{V}\,\diag(e^{\lambda_j\tau})\,\mat{V}^{-1}$,
and for large $\tau$ the rank-1 term
$e^{\lambda_{\max}\tau}\,\vect{v}_1\,\vect{w}_1^T$ dominates every
column of the matrix, where $\vect{v}_1$ and $\vect{w}_1^T$ are the
right and left eigenvectors for $\lambda_{\max}$.
\end{proposition}


%======================================================================
\section{Why the Step-by-Step SVD Fails}\label{sec:svd_failure}
%======================================================================

The current implementation accumulates the propagator in SVD form:
\[
  \mat{\Phi}(0\to k\cdot h) = \mat{U}_k\,\mat{\Sigma}_k\,\mat{V}_k^T,
\]
updated at each step via
$\mat{G} = \mat{\Phi}_{\text{step}}\,\mat{U}_k\,\mat{\Sigma}_k$,
then $\mat{G}=\hat{\mat{U}}\,\hat{\mat{\Sigma}}\,\hat{\mat{W}}^T$ (thin SVD).

\subsection{The mechanism of failure}
The update computes $\mat{G}=\mat{B}\,\diag(\mat{\Sigma}_k)$ where
$\mat{B}=\mat{\Phi}_{\text{step}}\mat{U}_k$.  When $\sigma_{\max}/\sigma_{\min}$
exceeds $\sim 10^{16}$ (machine precision in double float), the column
$\mat{G}_{:,j}=\mat{B}_{:,j}\cdot\sigma_j$ with $\sigma_j\ll\sigma_{\max}$
is polluted by roundoff from the large columns at the level
\[
  |\text{error in }\mat{G}_{:,j}| \;\sim\; \epsilon_{\text{mach}}\cdot
  \|\mat{B}_{:,1}\|\cdot\sigma_1,
  \qquad
  \frac{\text{relative error}}{\text{true value}} \;\sim\;
  \epsilon_{\text{mach}}\cdot\frac{\sigma_1}{\sigma_j}.
\]
Once $\sigma_1/\sigma_j > 10^{16}$, the $j$-th column of $\mat{G}$ is pure
noise, and the SVD of $\mat{G}$ returns a spurious $\sigma_j$.

\subsection{Diagnostic evidence}
For $\omega=0.5$, $\tbot=5$, $K=100$ steps, we tracked the step-by-step SVD
against the exact SVD of $\exp(\mat{A}\cdot k h)$:

\begin{table}[h]\centering
\caption{Step-by-step SVD vs.\ exact SVD,
$\omega=0.5$, $\tbot=5$, $K=100$.}\label{tab:svd_track}
\begin{tabular}{r|r|cc|cc}
\toprule
Step & $\tau$ & $\sigma_{\max}^{\text{step}}$ & $\sigma_{\max}^{\text{exact}}$
     & $\sigma_{\min}^{\text{step}}$ & $\sigma_{\min}^{\text{exact}}$ \\
\midrule
 10 & 0.50 & $9.95\times 10^{2}$  & $9.95\times 10^{2}$  & $1.01\times 10^{-3}$ & $1.01\times 10^{-3}$ \\
 20 & 1.00 & $9.72\times 10^{5}$  & $9.72\times 10^{5}$  & $1.03\times 10^{-6}$ & $1.03\times 10^{-6}$ \\
 30 & 1.50 & $9.50\times 10^{8}$  & $9.50\times 10^{8}$  & $1.05\times 10^{-9}$ & $7.86\times 10^{-10}$ \\
 40 & 2.00 & $9.28\times 10^{11}$ & $9.28\times 10^{11}$ & $1.08\times 10^{-12}$ & $3.92\times 10^{-7}$ \\
 50 & 2.50 & $9.07\times 10^{14}$ & $9.07\times 10^{14}$ & $1.10\times 10^{-15}$ & $1.04\times 10^{-6}$ \\
 54 & 2.70 & $1.42\times 10^{16}$ & $1.42\times 10^{16}$ & $\mathbf{7.0\times 10^{-17}}$ & $\mathbf{8.4\times 10^{-4}}$ \\
\bottomrule
\end{tabular}
\end{table}

At step 54 ($\tau\approx 2.7$), the step-by-step $\sigma_{\min}$ is
$7\times 10^{-17}$ (below machine epsilon) while the \emph{true}
$\sigma_{\min}$ is $8.4\times 10^{-4}$ --- a discrepancy of 13 orders of
magnitude.  The step-by-step SVD tracks $\sigma_{\max}$ perfectly (ratio
$\equiv 1.000000$) but catastrophically loses $\sigma_{\min}$ once the
condition number exceeds double precision.

\subsection{Why this is fundamental, not a bug}
The failure is not specific to SVD factorisation.  \emph{Any} algorithm
that accumulates the full $2N\times 2N$ propagator $\mat{\Phi}(0\to\tbot)$
will face the same problem: the propagator has condition number
$\cond\bigl(\mat{\Phi}\bigr) = \sigma_{\max}/\sigma_{\min} \approx
e^{2\lambda_{\max}\tbot}$, which exceeds double precision for
$\tbot \gtrsim 36/\lambda_{\max} \approx 2.6$ (for $\lambda_{\max}\approx 14$).

The non-normal transient growth (\Cref{sec:transient}) makes this worse:
even the ``decaying'' singular values rebound and become large, so one
cannot simply truncate them.  The freeze-and-keep approach (attempted in
a previous session) failed precisely because it assumed that small step-by-step
singular values corresponded to truly decaying physical modes --- but
the rebound phenomenon shows they do not.


%======================================================================
\section{Previous Failed Approaches}\label{sec:failures}
%======================================================================

\subsection{Attempt 1: Step-by-step SVD with freeze-and-keep}
\textbf{Idea.}  When $\sigma_j$ drops below a threshold, ``freeze'' that
mode: remove it from the live SVD and record its $\mat{V}^T$ row and
$q_{\text{scaled}}$ entry as fixed.

\textbf{Why it failed.}  The frozen $\mat{V}^T_{\text{gro}}$ subspace
(growing modes) was computed from the step-by-step SVD, which had already
diverged from the true SVD by the time of freezing.  Canonical cosines
between the step-by-step and exact growing-mode subspaces were as low as
0.073, meaning the subspaces were nearly orthogonal.  The BC solver,
which depends on accurate $\mat{V}^T_{\text{gro}}$, produced garbage.

\subsection{Attempt 2: Gram--Schmidt re-orthogonalisation}
\textbf{Idea.}  After freezing, re-orthogonalise the surviving modes
against the frozen growing-mode subspace using Gram--Schmidt, to prevent
mode leakage.

\textbf{Why it failed.}  The frozen $\mat{V}^T_{\text{gro}}$ was itself
inaccurate (see above).  Re-orthogonalising against an inaccurate subspace
does not recover accuracy---it merely projects onto the wrong complement.

\subsection{Attempt 3: Diffusion-domain handoff (current placeholder)}
\textbf{Idea.}  Detect when $\sigma_{\min}<\epsilon_{\text{mach}}$ and hand
off to a separate diffusion-domain solver.  The Magnus propagator at the
handoff point $\ttrans$ is accurate (all $\sigma_j>\epsilon_{\text{mach}}$),
and the diffusion solver handles the remaining depth $[\ttrans,\tbot]$.

\textbf{Status.}  The detection and handoff are implemented; the diffusion
solver itself is a placeholder (\texttt{NotImplementedError}).  The 6 failing
tests all trigger this handoff at
$\ttrans\approx\ln(1/\epsilon_{\text{mach}})/\lambda_{\max}\approx
36/\lambda_{\max}$.


%======================================================================
\section{Theoretical Framework}\label{sec:theory}
%======================================================================

\subsection{Magnus expansion: convergence and accuracy}
The Magnus expansion writes the solution of $\vect{I}'=\mat{A}(t)\vect{I}$
as $\vect{I}(t)=\exp\bigl(\mat{\Omega}(t)\bigr)\vect{I}(0)$, where
$\mat{\Omega}=\mat{\Omega}_1+\mat{\Omega}_2+\cdots$ is an infinite series
in nested integrals and commutators of $\mat{A}$~\cite{magnus1954,blanes2009}.

The first-order truncation $\mat{\Omega}_1=\int_0^h \mat{A}(s)\,ds$
approximated by the midpoint rule $\mat{\Omega}_1\approx h\,\mat{A}(h/2)$
gives the method used in \texttt{pydisort\_magnus}.  This is a
second-order method: the local error per step is $O(h^3)$, giving global
$O(h^2)$ convergence.

\textbf{Convergence radius.}  The Magnus series converges provided
$\int_0^h\|\mat{A}(s)\|_2\,ds < \pi$~\cite{blanes2009,moan2008}.  For the
DISORT matrix with $\|\mat{A}\|_2\approx 14.4$, this gives
$h<\pi/14.4\approx 0.22$.  With the typical step size
$h=\tbot/K \leq 0.32$ (for $K\geq 100$), we are near or within this radius.
The \emph{per-step} Magnus exponential is accurate.

\textbf{Key conclusion.}  The Magnus integrator itself is mathematically
viable for thick atmospheres.  The convergence radius limits the step size
$h$, not the total depth $\tbot$.  The challenge is purely in how the
per-step results are \emph{combined}.

\subsection{The Stamnes--Conklin substitution (eigendecomposition approach)}
In the original DISORT algorithm~\cite{stamnes1988}, the eigendecomposition
$\mat{A}=\mat{V}\,\diag(\pm\lambda_j)\,\mat{V}^{-1}$ is computed once per
layer.  The growing-mode coefficients are parametrised from $\tbot$:
$c_j^{\text{gro}}(t)=\xi_j^{\text{gro}}\,e^{-\lambda_j(\tbot-\tau)}$,
so that the exponential factor is always $\leq 1$.  This is the
Stamnes--Conklin scaling transformation, and it renders the boundary-condition
system unconditionally stable.

The Magnus approach was designed to \emph{avoid} per-layer eigendecomposition
(to handle continuously varying $\omega(\tau)$ and phase function).  We need
an analog of the Stamnes--Conklin trick that works with matrix exponentials
instead of eigendecompositions.

\subsection{The interaction principle and Redheffer star product}
\label{sec:star_product}

A slab from $\tau_a$ to $\tau_b$ is characterised by its
\emph{scattering matrix}: given radiation incident from both sides, what
comes out?  Formally, define the $N\times N$ reflection and transmission
operators:
\begin{align}
  \vect{I}^+(0) &= \mat{R}_\uparrow\,\vect{I}^-(0)
                  + \mat{T}_\uparrow\,\vect{I}^+(\tbot)
                  + \vect{s}_\uparrow, \label{eq:interaction_up} \\
  \vect{I}^-(\tbot) &= \mat{T}_\downarrow\,\vect{I}^-(0)
                      + \mat{R}_\downarrow\,\vect{I}^+(\tbot)
                      + \vect{s}_\downarrow. \label{eq:interaction_down}
\end{align}
Here $\mat{R}_\uparrow$ is the slab's reflection of downwelling radiation
(observed at the top), $\mat{T}_\downarrow$ is its downward transmission,
$\mat{R}_\downarrow$ is its reflection of upwelling radiation (observed at
the bottom), $\mat{T}_\uparrow$ is its upward transmission, and
$\vect{s}_\uparrow,\vect{s}_\downarrow$ are source terms from internal
scattering of the collimated beam.

When two adjacent slabs (labelled 1 on top, 2 below) are combined, the
\emph{Redheffer star product}~\cite{redheffer1962,grant1969} gives the
combined slab's operators:
\begin{align}
  \mat{R}_\uparrow^{1\star 2} &= \mat{R}_\uparrow^{(1)}
    + \mat{T}_\uparrow^{(1)}\,\mat{r}_\uparrow^{(2)}\,
      (\mat{I}-\mat{R}_\downarrow^{(1)}\mat{r}_\uparrow^{(2)})^{-1}\,
      \mat{T}_\downarrow^{(1)}, \label{eq:star_R} \\
  \mat{T}_\downarrow^{1\star 2} &= \mat{t}_\downarrow^{(2)}\,
      (\mat{I}-\mat{R}_\downarrow^{(1)}\mat{r}_\uparrow^{(2)})^{-1}\,
      \mat{T}_\downarrow^{(1)}, \label{eq:star_T}
\end{align}
and analogously for $\mat{R}_\downarrow^{1\star 2}$,
$\mat{T}_\uparrow^{1\star 2}$, and the source terms.

\begin{proposition}[Boundedness of star-product intermediates]
For a non-conservative atmosphere ($\omega<1$ everywhere), the operators
$\mat{R}_\uparrow,\mat{R}_\downarrow$ are bounded: their spectral norms
satisfy $\|\mat{R}\|_2<1$ (by energy conservation---some energy is absorbed
at each scattering event).  The transmission operators
$\mat{T}_\uparrow,\mat{T}_\downarrow$ decay as
$\|\mat{T}\|_2\sim e^{-\lambda_N\tau}$ (governed by the smallest eigenvalue
magnitude).  The resolvent
$(\mat{I}-\mat{R}_\downarrow^{(1)}\mat{r}_\uparrow^{(2)})^{-1}$ is
well-conditioned because
$\|\mat{R}_\downarrow^{(1)}\|\cdot\|\mat{r}_\uparrow^{(2)}\|<1$.
\end{proposition}

\textbf{Crucial insight.}  All intermediate $N\times N$ quantities in the
star product are $O(1)$.  No exponentially growing or decaying numbers
appear anywhere.  This is the fundamental difference from the propagator
accumulation approach.

\subsection{Extracting $\mat{R}$, $\mat{T}$ from the Magnus step}
From the $2N\times 2N$ propagator block decomposition
\[
  \mat{\Phi}_k = \begin{pmatrix} \mat{a} & \mat{b} \\ \mat{c} & \mat{d}
  \end{pmatrix},
  \qquad
  \vect{\delta}_k = \begin{pmatrix} \vect{\delta}^+ \\ \vect{\delta}^- \end{pmatrix},
\]
the single-step reflection, transmission, and source are:
\begin{equation}\label{eq:block_extract}
  \mat{r}_\uparrow = -\mat{a}^{-1}\mat{b}, \quad
  \mat{t}_\uparrow = \mat{a}^{-1}, \quad
  \mat{r}_\downarrow = \mat{c}\,\mat{a}^{-1}, \quad
  \mat{t}_\downarrow = \mat{d}-\mat{c}\,\mat{a}^{-1}\mat{b},
\end{equation}
\[
  \vect{s}_\uparrow^{(k)} = -\mat{a}^{-1}\vect{\delta}^+, \qquad
  \vect{s}_\downarrow^{(k)} = \vect{\delta}^- - \mat{c}\,\mat{a}^{-1}\vect{\delta}^+.
\]
For a single Magnus step of size $h$, the block $\mat{a}$ is well-conditioned:
\begin{equation}
  \cond(\mat{a}) \approx 1 + O(\lambda_{\max}h)
  \quad\text{(measured: } \cond(\mat{a})\leq 1.30 \text{ for } h=0.02\text{).}
\end{equation}
The extraction~\eqref{eq:block_extract} is therefore numerically stable for
small $h$.

\subsection{Connection to invariant imbedding and the matrix Riccati equation}

The interaction principle~\eqref{eq:interaction_up}--\eqref{eq:interaction_down}
is closely related to the \emph{invariant imbedding} formulation of
Bellman, Kalaba, and Prestrud~\cite{bellman1960}.  In that framework, the
reflection matrix $\mat{R}(\tau)$ for the slab $[0,\tau]$ satisfies the
matrix Riccati ODE
\begin{equation}\label{eq:riccati}
  \frac{d\mat{R}}{d\tau} = -\mat{\alpha}\mat{R} - \mat{R}\mat{\alpha}
                            - \mat{R}\mat{\beta}\mat{R} - \mat{\beta},
\end{equation}
where $\mat{\alpha}$ and $\mat{\beta}$ are the blocks of $\mat{A}$
in~\eqref{eq:A_block}.  The connection to the star product is:
the Riccati equation~\eqref{eq:riccati} is the \emph{infinitesimal
generator} of the star-product semigroup.  Specifically, the exact Riccati
step from $\tau$ to $\tau+h$ is given by the M\"obius transformation
\begin{equation}\label{eq:mobius}
  \mat{R}(\tau+h) = (\mat{a}\,\mat{R}(\tau)+\mat{b})\,
                     (\mat{c}\,\mat{R}(\tau)+\mat{d})^{-1}
\end{equation}
where $[\mat{a},\mat{b};\mat{c},\mat{d}]=\exp(\mat{A}\cdot h)$.

The Riccati formulation~\eqref{eq:riccati} is stable because $\mat{R}$ is
physically bounded ($\|\mat{R}\|_2<1$ for $\omega<1$).  No exponentially
growing quantities ever appear.


%======================================================================
\section{Is Magnus Integration Mathematically Viable?}\label{sec:viability}
%======================================================================

\textbf{Yes.}  The first-order Magnus integrator computes an accurate
matrix exponential $\exp(h\,\mat{A})$ at each step, provided
$h\|\mat{A}\|_2 < \pi$ (the convergence radius of the Magnus series).
This condition limits the \emph{step size}, not the total optical depth.
For any $\tbot$, one can choose $K$ large enough that $h=\tbot/K$ satisfies
the convergence criterion.

The per-step accuracy is $O(h^2)$ (midpoint rule), and the global accuracy
is $O(h^2)$ for $K$ steps.  This has been verified experimentally: with
1600 steps for $\tbot=32$ (the hardest test case), the star-product approach
achieves $0.25\%$ relative error, and the convergence ratio is
$3.97\approx 4.0 = 2^2$ (confirming second-order convergence).

The obstacle is not the Magnus integration itself but the propagator
accumulation.  The star product (\Cref{sec:star_product}) resolves this
completely.


%======================================================================
\section{Diagnostic Evidence}\label{sec:diagnostics}
%======================================================================

\subsection{Star-product convergence}
We implemented the star-product combination as a standalone diagnostic
and tested it against the pydisort eigendecomposition reference for all
6 failing test cases.

\begin{table}[h]\centering
\caption{Star-product convergence with step count (first-order Magnus).
``Ratio'' = error reduction when doubling $K$.}\label{tab:convergence}
\begin{tabular}{ll|rrr|rrr}
\toprule
Test & ($\tbot$, $\omega$)
  & \multicolumn{3}{c|}{Flux rel.\ error}
  & \multicolumn{3}{c}{Convergence ratio} \\
  & & $K\!=\!400$ & $K\!=\!800$ & $K\!=\!1600$
  & 400/800 & 800/1600 \\
\midrule
2c & (5, 0.5)     & $9.2\times 10^{-4}$ & $2.3\times 10^{-4}$ & $5.8\times 10^{-5}$ & 4.00 & 4.00 \\
1d & (32, 0.2)    & $4.0\times 10^{-2}$ & $1.0\times 10^{-2}$ & $2.6\times 10^{-3}$ & 3.89 & 3.97 \\
1f & (32, 0.99)   & $2.7\times 10^{-2}$ & $7.0\times 10^{-3}$ & $1.7\times 10^{-3}$ & 3.94 & 3.98 \\
\bottomrule
\end{tabular}
\end{table}

The convergence ratio approaches 4.0 in all cases, confirming the $O(h^2)$
rate expected from the midpoint-rule Magnus.  \emph{No numerical instability
is observed for any $\tbot$.}

\subsection{Boundedness of intermediates}
At every step of the star-product accumulation, we verified:
\begin{itemize}[nosep]
\item $\cond(\mat{a}_k)\leq 1.30$ (block of each Magnus step),
\item $\cond(\mat{I}-\mat{r}_\uparrow^{(k)}\mat{R}_\downarrow)\leq 1.06$
  (resolvent in star product),
\item $\|\mat{R}_\uparrow\|_2 \leq 0.94$ (reflection always sub-unitary),
\item $\|\mat{T}\|_2 \leq 1.0$ (transmission bounded).
\end{itemize}
These confirm unconditional numerical stability.

\subsection{Handling of boundary conditions}
The star-product approach correctly handles all boundary condition types.
With $K=400$ steps, the $m=0$ mode matches pydisort to relative errors
$\leq 10^{-5}$:
\begin{center}
\begin{tabular}{lrr}
\toprule
Test case & Star-product flux & Rel.\ error \\
\midrule
4a: $b_{\text{neg}}=1/\pi$, no beam   & $1.076167\times 10^{-1}$ & $2.6\times 10^{-16}$ \\
4b: beam + $b_{\text{pos}}=0.5$, HG   & $9.936779\times 10^{-1}$ & $1.5\times 10^{-7}$  \\
4c: both BCs, isotropic                & $6.088967\times 10^{-1}$ & $6.5\times 10^{-6}$  \\
5a: BDRF $\rho=0.1$, isotropic         & $6.473253\times 10^{-2}$ & $7.4\times 10^{-7}$  \\
5b: BDRF $\rho=0.5$, HG               & $8.631945\times 10^{-2}$ & $1.8\times 10^{-6}$  \\
\bottomrule
\end{tabular}
\end{center}


%======================================================================
\section{Proposed Approaches}\label{sec:proposals}
%======================================================================

\subsection{Approach 1: Star-product combination of Magnus steps
(primary recommendation)}\label{sec:approach1}

\textbf{Core idea.}  Replace the step-by-step SVD accumulation of the
$2N\times 2N$ propagator with the Redheffer star-product accumulation of the
$N\times N$ reflection, transmission, and source operators.

\textbf{Per step.}  Compute the Magnus exponential $\mat{\Phi}_k$ and
particular increment $\vect{\delta}_k$ exactly as now (augmented
\texttt{expm}).  Extract $\mat{r}_\uparrow$, $\mat{t}_\uparrow$,
$\mat{r}_\downarrow$, $\mat{t}_\downarrow$, $\vect{s}_\uparrow$,
$\vect{s}_\downarrow$ via~\eqref{eq:block_extract}.

\textbf{Accumulation.}  Combine each new step's operators with the
accumulated operators using the star-product
formulas~\eqref{eq:star_R}--\eqref{eq:star_T} and analogous source
formulas (derived in the diagnostics).

\textbf{BC solve.}  From the accumulated $(\mat{R}_\uparrow, \mat{T}_\uparrow,
\mat{T}_\downarrow, \mat{R}_\downarrow, \vect{s}_\uparrow,
\vect{s}_\downarrow)$, solve the boundary-condition system:
\begin{align}
  (\mat{I}-\mat{R}_{\text{surf}}\mat{R}_\downarrow)\,\vect{I}^+(\tbot)
    &= \mat{R}_{\text{surf}}(\mat{T}_\downarrow\,\vect{b}_{\text{neg}}
       + \vect{s}_\downarrow) + \vect{b}_{\text{pos}}^{\text{eff}}, \\
  \vect{I}^+(0) &= \mat{R}_\uparrow\,\vect{b}_{\text{neg}}
                  + \mat{T}_\uparrow\,\vect{I}^+(\tbot)
                  + \vect{s}_\uparrow.
\end{align}
This is an $N\times N$ linear system (half the size of the current
$2N\times 2N$ Stamnes--Conklin system) followed by a matrix-vector product.

\textbf{Advantages.}
\begin{itemize}[nosep]
\item Unconditionally stable: all intermediates are $O(1)$.
\item Clean $O(h^2)$ convergence, verified for $\tbot$ up to 32.
\item Naturally handles BDRF, $b_{\text{pos}}$, $b_{\text{neg}}$.
\item The per-step Magnus computation is unchanged.
\item Cost: $O(K\cdot N^3)$ --- same as the current approach.
\item Memory: $O(N^2)$ --- only the current accumulated operators are stored.
\end{itemize}

\textbf{Disadvantages.}
\begin{itemize}[nosep]
\item Requires rewriting \texttt{\_compute\_magnus\_propagator} (accumulation
  loop) and \texttt{\_solve\_bc\_magnus} (BC solve).  The interface to
  \texttt{pydisort\_magnus} is unchanged.
\item The number of steps $K$ must be large enough that the per-step block
  $\mat{a}$ is well-conditioned:
  $h\cdot\lambda_{\max}\lesssim 1$--$2$ in practice.
  For $\tbot=32$: $K\gtrsim 32\cdot 14/2 = 224$.
\end{itemize}


\subsection{Approach 2: Bidirectional Magnus propagation with
mode-partitioned interface matching}\label{sec:approach2}

\textbf{Core idea.}  Instead of propagating all $2N$ modes in one direction
(forward), split into two half-rank sweeps:
\begin{enumerate}[nosep]
\item \textbf{Forward sweep} ($\tau=0\to\ttrans$): propagate $N$
  fundamental solutions that are initially aligned with the
  \emph{decaying} eigenvectors of $\mat{A}$.  These solutions shrink
  (on average) going forward, so the sweep is well-conditioned.
\item \textbf{Backward sweep} ($\tau=\tbot\to\ttrans$): propagate $N$
  fundamental solutions aligned with the \emph{growing} eigenvectors.
  Going backward, these become decaying, so this sweep is also
  well-conditioned.
\item \textbf{Match at $\ttrans$}: the $2N$ columns from both sweeps form
  a complete basis at the interface.  Apply boundary conditions to
  determine the coefficients.
\end{enumerate}

\textbf{Continuous orthogonalisation.}  In each sweep, the non-normality
of $\mat{A}$ causes ``leakage'' of roundoff into the unwanted mode
directions.  At each step, we orthogonalise the $N$ columns and project
away the component in the unwanted eigenspace.  Concretely:
\begin{enumerate}[nosep]
\item Propagate: $\mat{Y}_{k+1} = \mat{\Phi}_k\,\mat{Y}_k$
  \quad ($\mat{Y}\in\R^{2N\times N}$).
\item QR: $\mat{Y}_{k+1} = \mat{Q}_{k+1}\,\mat{R}_{k+1}$.
\item Project: remove components along the $N$ unwanted directions
  (using the eigenvectors of $\mat{A}(\tau_{k+1})$ or, for slowly varying
  $\mat{A}$, of a recent $\mat{A}$).
\item Store $\mat{R}_{k+1}$ factors separately; maintain
  $\log\det(\mat{R}_{\text{total}})$ for amplitude tracking.
\end{enumerate}

\textbf{Required eigendecomposition.}  The projection in step~3 requires
knowing the growing/decaying eigenvectors at each step (or at least
periodically).  For \emph{constant} $\mat{A}$ (all currently failing tests),
a single eigendecomposition suffices.  For variable $\mat{A}(\tau)$, the
eigenvectors must be tracked, adding $O(N^3)$ per step.

\textbf{Transient growth bound.}  The roundoff contamination from the
unwanted modes in each step is $O(\epsilon_{\text{mach}}\cdot\cond(\mat{V})^2
\cdot e^{\lambda_{\max}h})$, where $\cond(\mat{V})\leq 3.7$ for the DISORT
$\mat{A}$.  After projection, this contamination is removed.  The
\emph{accumulated} error after $K$ steps is
$O(K\cdot\epsilon_{\text{mach}}\cdot\cond(\mat{V})^2)$, which remains
small ($\lesssim 10^{-12}$ for $K=1600$).

\textbf{Advantages.}
\begin{itemize}[nosep]
\item Conceptually closer to the current code structure (still propagates
  fundamental solutions, not R/T operators).
\item The interface matching is a $2N\times 2N$ linear system, similar to
  the current BC solver.
\item For constant $\mat{A}$: only one eigendecomposition needed, and the
  projection is exact.
\end{itemize}

\textbf{Disadvantages.}
\begin{itemize}[nosep]
\item Requires eigendecomposition (at least once), partially defeating the
  purpose of the Magnus approach (which was to avoid eigendecomposition).
  For variable $\mat{A}(\tau)$, eigendecompositions are needed periodically.
\item More complex implementation: two sweeps, QR + projection at each step,
  $\mat{R}$-factor accumulation, interface matching.
\item The projection step is specific to the $\pm\lambda$ eigenvalue pairing
  of the DISORT $\mat{A}$; it may not generalise to other ODE structures.
\item Less tested than Approach~1 (no diagnostic prototype exists yet).
\end{itemize}


%======================================================================
\section{Comparison of Approaches}\label{sec:comparison}
%======================================================================

\begin{table}[h]\centering
\caption{Comparison of the two proposed approaches.}\label{tab:comparison}
\begin{tabular}{lcc}
\toprule
Property & Approach 1 (Star product) & Approach 2 (Bidirectional) \\
\midrule
Stability           & Unconditional       & Conditional (needs projection) \\
Eigendecomp needed? & No                  & Yes (at least once) \\
Cost per step       & $O(N^3)$            & $O(N^3)$ \\
BC system size      & $N\times N$         & $2N\times 2N$ \\
Implementation      & Moderate rewrite    & Major rewrite \\
Tested?             & Yes (diagnostics)   & No \\
Variable $\omega(\tau)$? & Native          & Needs periodic eigendecomp \\
\bottomrule
\end{tabular}
\end{table}


%======================================================================
\section{Conclusion and Recommendation}
%======================================================================

The Magnus integrator is \textbf{mathematically viable} for thick atmospheres.
The per-step matrix exponential is accurate whenever $h\|\mat{A}\|_2<\pi$,
which is easily achievable by choosing a sufficient number of steps.  The
obstacle is purely the \emph{propagator accumulation strategy}.

\textbf{Primary recommendation: Approach 1 (star product).}  This approach
is proven to work (diagnostic evidence for all failing test cases), is
unconditionally stable, requires no eigendecomposition, naturally handles
all boundary condition types, and has the same computational cost as the
current approach.  The convergence is $O(h^2)$ with a clean ratio of 4.0.

Approach~2 (bidirectional propagation) is a viable alternative that stays
closer to the current code structure, but it requires eigendecomposition
and has not been prototyped.  It may be preferable if the user wishes to
preserve the current interface between the propagator and BC solver.

The current diffusion-domain handoff (Attempt~3 / placeholder) is
\emph{not recommended} as the primary strategy, because it requires a
separate solver for $[\ttrans,\tbot]$ that would itself need to solve
the same stability problem.  However, it could be used as a
\emph{complement} to Approach~1 for extremely thick atmospheres where
$K$ would otherwise be impractically large.


%======================================================================
\begin{thebibliography}{99}
\bibitem{magnus1954}
W.~Magnus, ``On the exponential solution of differential equations for a
linear operator,'' \emph{Comm.\ Pure Appl.\ Math.}, 7:649--673, 1954.

\bibitem{blanes2009}
S.~Blanes, F.~Casas, J.A.~Oteo, J.~Ros, ``The Magnus expansion and some
of its applications,'' \emph{Physics Reports}, 470(5--6):151--238, 2009.

\bibitem{moan2008}
P.C.~Moan and J.~Niesen, ``Convergence of the Magnus series,''
\emph{Found.\ Comput.\ Math.}, 8(3):291--301, 2008.

\bibitem{stamnes1988}
K.~Stamnes, S.-C.~Tsay, W.~Wiscombe, K.~Jayaweera, ``Numerically stable
algorithm for discrete-ordinate-method radiative transfer in multiple
scattering and emitting layered media,'' \emph{Appl.\ Opt.}, 27(12):
2502--2509, 1988.

\bibitem{stamnes1984}
K.~Stamnes and P.~Conklin, ``A new multi-layer discrete ordinate approach to
radiative transfer in vertically inhomogeneous atmospheres,'' \emph{J.\ Quant.\
Spectrosc.\ Radiat.\ Transfer}, 31(3):273--282, 1984.

\bibitem{kreiss1962}
H.-O.~Kreiss, ``\"Uber die Stabilit\"atsdefinition f\"ur
Differenzengleichungen die partielle Differentialgleichungen
approximieren,'' \emph{BIT}, 2:153--181, 1962.

\bibitem{redheffer1962}
R.~Redheffer, ``On the relation of transmission-line theory to scattering and
transfer,'' \emph{J.\ Math.\ Phys.}, 41:1--41, 1962.

\bibitem{grant1969}
I.P.~Grant and G.E.~Hunt, ``Discrete space theory of radiative transfer.
I.~Fundamentals,'' \emph{Proc.\ R.\ Soc.\ Lond.\ A}, 313(1513):183--197, 1969.

\bibitem{bellman1960}
R.~Bellman, R.~Kalaba, and M.~Prestrud, \emph{Invariant Imbedding and
Radiative Transfer in Slabs of Finite Thickness}, American Elsevier, 1960.

\bibitem{plass1973}
G.N.~Plass, G.W.~Kattawar, and F.E.~Catchings, ``Matrix operator theory of
radiative transfer.\ 1: Rayleigh scattering,'' \emph{Appl.\ Opt.}, 12(2):
314--329, 1973.

\bibitem{lacis1998}
A.~Lacis and V.~Oinas, ``A description of the correlated $k$ distribution
method for modeling nongray gaseous absorption,'' \emph{J.\ Geophys.\ Res.},
96(D5):9027--9063, 1991.
\end{thebibliography}

\end{document}
